% Generated by tools/edn_latex.py from reports/edn.md.
% Do not edit: change reports/edn.md and run `make edn`.
\ednentry{
  id = {EDN-62},
  title = {\texorpdfstring{\texttt{lgbm-\allowbreak{}basic}}{lgbm-basic} is the candidate, and SC-04's miss on cars over 20 years is accepted rather than worked around},
  date = {2026-09-30},
  milestone = {M3: Quality Assurance},
  topic = {Modelling Approach, Model Evaluation (the NFR-01 gate)},
  participants = {@lukas2510, decided by the repository owner},
  decision = {Two decisions the measured criteria forced, recorded together because the same measurement produced both. \textbf{First:} \texttt{lgbm-\allowbreak{}basic} is the candidate the model card and the report put forward, not \texttt{lgbm-\allowbreak{}extended}, although \texttt{lgbm-\allowbreak{}extended} has the lower pooled MdAPE (6.32\,\% against 6.83\,\%). \texttt{lgbm-\allowbreak{}extended} fails SC-06 at 1.584x and \texttt{lgbm-\allowbreak{}basic} passes at 1.208x, and on a request carrying only the ten fields FR-01 requires -- which the contract lets any caller send -- \texttt{lgbm-\allowbreak{}basic} answers at 8.25\,\% against \texttt{lgbm-\allowbreak{}extended}'s 10.01\,\%. The candidate is a documented judgement, not a computed one: \texttt{metrics.\allowbreak{}json}'s \texttt{best\_\allowbreak{}variant} stays the lowest MdAPE overall, because that field means what it says, and \texttt{deployable\_\allowbreak{}variant} stays null while any criterion blocks. \textbf{Second:} SC-04's failure on \texttt{age\_\allowbreak{}bucket=\allowbreak{}over 20} (17.12\,\% for the candidate against a 15\,\% bound, over 784 test rows) is recorded as a missed criterion. The threshold is not moved, no upper age bound is added to the scope, and the segment is not dropped from the breakdown.},
  rationale = {\emph{Alternatives considered:}
\begin{itemize}[leftmargin=*, topsep=2pt]
\item \textbf{Option A (chosen): put \texttt{lgbm-\allowbreak{}basic} forward, and accept SC-04's miss with the reason recorded.}
\newline Pros: it is the only reading under which SC-06 does any work. A criterion exists to change a decision, and this is the decision it changes. The candidate meets every criterion that is measurable except SC-04, which no variant meets, so the choice costs nothing in criteria met and 0.51 pp of pooled MdAPE on a fully described car. It is also the cheaper model by every other measure: 16 columns against 162, 17.8 s of fit against 43 s, 6.3 MB of payload against 6.8 MB, and no 90-second \texttt{features} stage for the extended matrix. Accepting SC-04's miss keeps the bar where EDN-06 set it and turns the failure into the finding it is.
\newline Cons: the report has to explain why the model with the better headline number is not the one put forward, which is more words than ``we picked the best''. And the candidate is now a statement in a document rather than a field in an artefact, so a reader of \texttt{metrics.\allowbreak{}json} alone does not see it.
\item \textbf{Option B: keep \texttt{lgbm-\allowbreak{}extended} as the candidate and note SC-06 as a known limitation.}
\newline Pros: the lowest MdAPE is the number a reader looks for first, and the criteria are advisory until a release actually happens.
\newline Cons: it makes SC-06 decorative. The criterion was added by EDN-15 precisely to catch a model whose accuracy depends on optional inputs, it caught one, and overriding it the first time it fires teaches everyone reading the report that the gate is a formality. It would also put forward the \emph{worse} model for the request FR-01 guarantees, by 1.76 pp.
\item \textbf{Option C: serve \texttt{lgbm-\allowbreak{}extended} when the request is complete and \texttt{lgbm-\allowbreak{}basic} when it is not.}
\newline Pros: strictly the best point estimate for every request, and the 0.51 pp is not thrown away.
\newline Cons: two models in the image, two sets of SHAP baselines, two rows in every metrics table, and a routing rule that has to be explained to a user whose estimate changes when they add an optional field. It is also not this ticket's decision to take: it is an API design question, and FR-02's answer may turn out to be an interval rather than a second model. Left open for the deployment milestone rather than ruled out.
\item \textbf{Option D (for SC-04): add an upper age bound to the training scope, the way EDN-04 bounds price.}
\newline Pros: the criterion would pass, and there is precedent in the same document.
\newline Cons: it would pass by removing the rows it fails on, which is the shape of a criterion adjusted to its result. The cars over 20 years old are 784 of 19,665 test rows and they are cars people ask about; declaring them out of scope is a product decision about who the component serves, and it cannot be justified by a criterion the model missed.
\item \textbf{Option E (for SC-04): raise the threshold to 20\,\%, or report the criterion on a subset.}
\newline Pros: it would make the gate green.
\newline Cons: it is the same objection without the honesty of a stated scope change. SC-04's 15\,\% came from a reference run and a deliberate margin, and nothing has been learned about the \emph{bar} -- only about the model.
\end{itemize}
\par
\emph{Why:} The two halves of this entry rest on one measurement, and the measurement is what makes the choice defensible rather than tasteful. Both LightGBM variants were swept over every optional field on the real snapshot: each of \texttt{lgbm-\allowbreak{}extended}'s 23 optional fields costs at most 9.2\,\% of its MdAPE on its own, but all 23 at once cost 58.4\,\%, and that all-at-once scenario is not a hypothetical -- it is exactly a request carrying only the required fields, which FR-01 promises to answer. Comparing the two variants \emph{on that request} is comparing them on a request the product has to serve, and there the ladder's conclusion reverses. So the honest statement of what ladder step 4 measured is conditional: the extended features are worth 0.51 pp when the caller fills them in and cost 1.76 pp when the caller does not, and for a component that lets a caller omit 23 of its 33 inputs the second half is the one that decides.
\par
SC-04 is the same measurement read the other way. The problem specification named two possible remedies for the over-20-years segment, ``the extended features or a scope change for classic cars'', and the extended features have now been measured against it: 17.78\,\% against the basic set's 17.12\,\%, so they make that segment marginally \emph{worse} while improving the pooled figure. That rules the first remedy out on evidence, which is a finding about the features rather than an argument about the bar, and it is worth recording as such. The second remedy is a product decision nobody has taken, so the criterion is reported as missed. Nothing is unblocked by softening it in any case: SC-05 has no measurement until the UC2 conformal intervals exist, so NFR-01's gate blocks whatever SC-04 says, and a threshold edited under those circumstances would have no effect except on the report's credibility.},
  ai = {Information seeking, Alternative generation, Alternative assessment, Recommendation},
  response = {Accepted},
  assessment = {The finding is AI's, and it was not the one anybody was looking for. Two sessions measured SC-06 independently -- the implementer of issue \#39 and the reviewer of issue \#37 -- and agreed to four decimal places on every variant, which is why the numbers were trusted. The implementer then noticed what the all-at-once scenario \emph{is}, namely a required-fields-only request, and that comparing the two variants on it inverts the ladder's ranking; that step is what turned a criterion result into a decision. AI recommended option A for both halves and raised options C, D and E as things it would not decide on its own, which is the right boundary: which model the API serves and which cars the product covers are not conclusions to draw from a metrics file. Two of its numbers were corrected before they were used: an expected single-field failure on \texttt{has\_\allowbreak{}full\_\allowbreak{}service\_\allowbreak{}history} (predicted at 1.399x) measured at 1.002x on real data and was reported as not reproducing rather than as confirmed, and \texttt{b1}'s figures were re-measured after it turned out the available models had been fitted at a stale \texttt{min\_\allowbreak{}category\_\allowbreak{}rows}.},
  aievidence = {Claude Code sessions on 2026-09-30 implementing and reviewing issues \#37 and \#39. The masking sweep over 19,665 real test rows produced 1.208x for \texttt{lgbm-\allowbreak{}basic} and 1.584x for \texttt{lgbm-\allowbreak{}extended}, reproduced independently by both sessions, and \texttt{reports/\allowbreak{}metrics/\allowbreak{}masked-\allowbreak{}inputs.\allowbreak{}csv} carries the per-field table behind those two numbers. The owner was given options A to E with these measurements and chose A.},
  evidence = {\href{https://github.com/mlops-2627q1-mds-upc/MLOPS_Recommenditos/blob/main/docs/docs/model-card.md}{model card}, Results and the experiment ladder; \texttt{reports/\allowbreak{}metrics/\allowbreak{}masked-\allowbreak{}inputs.\allowbreak{}csv} and \texttt{reports/\allowbreak{}metrics/\allowbreak{}segments.\allowbreak{}csv}; \href{https://github.com/mlops-2627q1-mds-upc/MLOPS_Recommenditos/blob/main/docs/docs/problem-spec.md}{problem specification} section 8; \hyperref[EDN-06]{EDN-06}, \hyperref[EDN-15]{EDN-15}, \hyperref[EDN-59]{EDN-59}; \href{https://github.com/mlops-2627q1-mds-upc/MLOPS_Recommenditos/issues/39}{issue \#39}, \href{https://github.com/mlops-2627q1-mds-upc/MLOPS_Recommenditos/pull/65}{PR \#65}.},
}
